\thispagestyle{empty}
\pdfbookmark[0]{Ausgabe und Quellen}{edition-and-sources}
\begin{center}
\textsc{Carl Friedrich Gauss}\par\medskip
{\Large Theorie der Anziehung homogener elliptischer sphäroidischer Körper\par}
\vspace{0pt}{\fontsize{10.75}{10.7}\selectfont Über ein neues allgemeines Grundgesetz der Mechanik\par}
\vspace{0pt}{\fontsize{10.75}{10.7}\selectfont Allgemeine Grundsätze der Theorie der Gestalt von Flüssigkeiten im Gleichgewicht\par}
\vspace{0pt}{\fontsize{10.75}{10.7}\selectfont Die Intensität der erdmagnetischen Kraft in absoluten Einheiten\par}
\vspace{0pt}{\fontsize{10.75}{10.7}\selectfont Allgemeine Theorie des Erdmagnetismus\par}
\vspace{0pt}{\fontsize{10.75}{10.7}\selectfont Allgemeine Lehrsätze über Anziehungs- und Abstoßungskräfte\par}
\vspace{0pt}{\fontsize{10.75}{10.7}\selectfont Dioptrische Untersuchungen\par}
\vspace{0pt}{\fontsize{10.75}{10.7}\selectfont Anzeigen eigener Abhandlungen\par}
\vspace{0pt}{\fontsize{10.75}{10.7}\selectfont Einleitung zu den Beobachtungen des Magnetischen Vereins von 1836\par}
\vspace{0pt}{\fontsize{10.75}{10.7}\selectfont Ein neues Hilfsmittel für die magnetischen Beobachtungen\par}
\vspace{0pt}{\fontsize{10.75}{10.7}\selectfont Ein neues Instrument zur Messung von Intensitätsänderungen\par}
\vspace{0pt}{\fontsize{10.75}{10.7}\selectfont Anleitung zur Bestimmung der Schwingungsdauer einer Magnetnadel\par}
\vspace{0pt}{\fontsize{10.75}{10.7}\selectfont Über ein Mittel, die Beobachtung von Ablenkungen zu erleichtern\par}
\vspace{0pt}{\fontsize{10.75}{10.7}\selectfont Zur Bestimmung der Konstanten des Bifilarmagnetometers\par}
\vspace{0pt}{\fontsize{10.75}{10.7}\selectfont Vorschriften zur Bestimmung der magnetischen Wirkung, die ein Magnetstab in der Ferne ausübt\par}
\vspace{0pt}{\fontsize{10.75}{10.7}\selectfont Über die Anwendung des Magnetometers zur Bestimmung der absoluten Deklination\par}
\vspace{0pt}{\fontsize{10.75}{10.7}\selectfont Beobachtungen der magnetischen Inklination in Göttingen\par}
\vspace{0pt}{\fontsize{10.75}{10.7}\selectfont Fundamentalgleichungen für die Bewegung schwerer Körper auf der rotierenden Erde\par}
\vspace{0pt}{\fontsize{10.75}{10.7}\selectfont Über achromatische Doppelobjektive\par}
\vspace{0pt}{\fontsize{10.75}{10.7}\selectfont Brief von Gauss an Brandes\par}
\vspace{0pt}{\fontsize{10.75}{10.7}\selectfont Berichtigung der Schneiden einer Waage\par}
\vspace{0pt}{\fontsize{10.75}{10.7}\selectfont Physikalische Beobachtungen\par}
\vspace{0pt}{\fontsize{10.75}{10.7}\selectfont Nordlicht am 7. Januar 1831\par}
\vspace{0pt}{\fontsize{10.75}{10.7}\selectfont Magnetisches Observatorium in Göttingen (August 1834)\par}
\vspace{0pt}{\fontsize{10.75}{10.7}\selectfont Beobachtungen der magnetischen Variation in Göttingen und Leipzig\par}
\vspace{0pt}{\fontsize{10.75}{10.7}\selectfont Magnetische Beobachtungen (März 1835)\par}
\vspace{0pt}{\fontsize{10.75}{10.7}\selectfont Magnetnadel in Kopenhagen und Mailand\par}
\vspace{0pt}{\fontsize{10.75}{10.7}\selectfont Brief von Gauss an Schumacher (1836)\par}
\vspace{0pt}{\fontsize{10.75}{10.7}\selectfont Verfahren bei den Beobachtungsterminen\par}
\vspace{0pt}{\fontsize{10.75}{10.7}\selectfont Dreijährige tägliche Deklinationsbeobachtungen\par}
\vspace{0pt}{\fontsize{10.75}{10.7}\selectfont Terminszeichnungen und Beobachtungszahlen von 1836\par}
\vspace{0pt}{\fontsize{10.75}{10.7}\selectfont Terminszeichnungen und Beobachtungszahlen von 1837\par}
\vspace{0pt}{\fontsize{10.75}{10.7}\selectfont Zu Wilkes’ Beobachtungen in der Nähe des magnetischen Südpols\par}
\vspace{0pt}{\fontsize{10.75}{10.7}\selectfont Besprechung von Benzenbergs Über die Daltonsche Theorie\par}
\vspace{0pt}{\fontsize{10.75}{10.7}\selectfont Besprechung von Fischers Anleitung zur Herstellung künstlicher Magnete\par}
\vspace{0pt}{\fontsize{10.75}{10.7}\selectfont Besprechung der Ergebnisse des Magnetischen Vereins für 1836\par}
\vspace{0pt}{\fontsize{10.75}{10.7}\selectfont Das mathematisch-physikalische Institut in Marburg\par}
\vspace{0pt}{\fontsize{10.75}{10.7}\selectfont Zur mathematischen Theorie der elektrodynamischen Wirkungen, §§ 1–21, Kugelfunktionen, Zum Gebrauch des Komparators, Wirkung eines leuchtenden Punkts, Scherings Bemerkungen, Historisches Inhaltsverzeichnis und Druckvermerk\par}
\smallskip
KI-Übersetzung aus dem Lateinischen und Übertragung in modernes Deutsch mit KI-erstelltem Satz.\\Deutsche Ausgabe
\end{center}
\clearpage
\begingroup\small
\smallskip
\noindent Die vollständige Abhandlung umfasst alle vierzehn Artikel, die historische Anmerkung und den Zusatz. Die Titelblätter des Bandes sind enthalten. Grundlage ist der lateinische Text in Carl Friedrich Gauss, \textit{Werke}, Band V, herausgegeben von der Königlichen Gesellschaft der Wissenschaften zu Göttingen (1867), S.~3--22, mit den historischen Titelblättern. Die Titelvignette ist dem von den \href{https://archive.org/details/Werkecarlf05gausrich/page/n8/mode/1up}{University of California Libraries} digitalisierten zweiten Abdruck (Göttingen, 1877) entnommen.
\par\medskip
\noindent \href{https://github.com/KokunoYumeto/gauss-collected-works/releases/tag/gauss-original-language-2026-10-24}{Originalsprachige Ausgabe}
\par\medskip
\noindent\textbf{Historische Fassungen und verwandte Darstellung}\par\smallskip
\noindent Zum Vergleich dienten die deutsche Fassung von 1813 in der \textit{Monatlichen Correspondenz zur Beförderung der Erd- und Himmels-Kunde}, Bd.~28, S.~37--41, und die deutsche Fassung in Albert Wangerin (Hrsg.), \textit{Ueber die Anziehung homogener Ellipsoide} (Leipzig, 1890), S.~50--74. Arthur Cayleys \textit{On Gauss’s Method for the Attraction of Ellipsoids} (1857), nachgedruckt in den \textit{Collected Mathematical Papers}, Bd.~III (1890), S.~25--28, ist eine eigenständige Darstellung.
\par\medskip
\noindent \href{https://www.digitale-sammlungen.de/de/view/bsb10538617?page=43}{1813}\quad\href{https://archive.org/details/ueberdieanziehu00wanggoog}{1890}\quad\href{https://rcin.org.pl/dlibra/publication/184416/edition/170906/on-gauss-method-for-the-attraction-of-ellipsoids-cayley-arthur-1821-1895?language=en}{Cayley}
\par\medskip
\noindent\textbf{Mechanik}\par\smallskip
\noindent Die vollständige Abhandlung von 1829 mit beiden Anmerkungen von Gauß ist eine KI-Übertragung des deutschen Originals in modernes Deutsch. Sie erschien im \textit{Journal für die reine und angewandte Mathematik}, Bd.~4 (1829), S.~232--235. Grundlage ist der Text in den \textit{Werken}, Band V, S.~25--28, mit dem Titelblatt. Die englische Übersetzung von D. H. Delphenich ist zum Vergleich als einzelnes PDF und in seinen \textit{Selected Papers on Gauss’s Principle of Least Constraint} (Neo-Classical Press, 2020) zugänglich.
\par\nobreak\medskip
\noindent \href{https://doi.org/10.1515/crll.1829.4.232}{Original von 1829}\quad\href{https://www.neo-classical-physics.info/uploads/3/4/3/6/34363841/gauss_-_princ_least_constraint.pdf}{Delphenich}\quad\href{https://www.neo-classical-physics.info/uploads/3/4/3/6/34363841/selecta_on_least_constraint.pdf#page=95}{Sammlung von 2020}
\par\medskip
\Needspace{4\baselineskip}
\noindent\textbf{Flüssigkeitsgleichgewicht}\par\smallskip
\noindent Die vollständige lateinische Abhandlung mit allen 34 Artikeln und den vier Anmerkungen von Gauß ist KI-übersetzt. Grundlage sind die \textit{Werke}, Band V, S.~31--77, mit dem Titelblatt von 1830. Rudolf Heinrich Webers deutsche Übersetzung von 1903, herausgegeben von Heinrich Weber in Ostwalds Reihe (Nr.~135), ist eine eigenständige historische Fassung.
\par\smallskip
\noindent \href{https://archive.org/details/allgemeinegrund00gausgoog}{Webers deutsche Ausgabe (1903)}
\par\endgroup
\clearpage
